% Options for packages loaded elsewhere
\PassOptionsToPackage{unicode}{hyperref}
\PassOptionsToPackage{hyphens}{url}
%
\documentclass[
]{article}
\usepackage{amsmath,amssymb}
\usepackage{iftex}
\ifPDFTeX
  \usepackage[T1]{fontenc}
  \usepackage[utf8]{inputenc}
  \usepackage{textcomp} % provide euro and other symbols
\else % if luatex or xetex
  \usepackage{unicode-math} % this also loads fontspec
  \defaultfontfeatures{Scale=MatchLowercase}
  \defaultfontfeatures[\rmfamily]{Ligatures=TeX,Scale=1}
\fi
\usepackage{lmodern}
\ifPDFTeX\else
  % xetex/luatex font selection
\fi
% Use upquote if available, for straight quotes in verbatim environments
\IfFileExists{upquote.sty}{\usepackage{upquote}}{}
\IfFileExists{microtype.sty}{% use microtype if available
  \usepackage[]{microtype}
  \UseMicrotypeSet[protrusion]{basicmath} % disable protrusion for tt fonts
}{}
\makeatletter
\@ifundefined{KOMAClassName}{% if non-KOMA class
  \IfFileExists{parskip.sty}{%
    \usepackage{parskip}
  }{% else
    \setlength{\parindent}{0pt}
    \setlength{\parskip}{6pt plus 2pt minus 1pt}}
}{% if KOMA class
  \KOMAoptions{parskip=half}}
\makeatother
\usepackage{xcolor}
\usepackage[margin=1in]{geometry}
\usepackage{graphicx}
\makeatletter
\newsavebox\pandoc@box
\newcommand*\pandocbounded[1]{% scales image to fit in text height/width
  \sbox\pandoc@box{#1}%
  \Gscale@div\@tempa{\textheight}{\dimexpr\ht\pandoc@box+\dp\pandoc@box\relax}%
  \Gscale@div\@tempb{\linewidth}{\wd\pandoc@box}%
  \ifdim\@tempb\p@<\@tempa\p@\let\@tempa\@tempb\fi% select the smaller of both
  \ifdim\@tempa\p@<\p@\scalebox{\@tempa}{\usebox\pandoc@box}%
  \else\usebox{\pandoc@box}%
  \fi%
}
% Set default figure placement to htbp
\def\fps@figure{htbp}
\makeatother
\setlength{\emergencystretch}{3em} % prevent overfull lines
\providecommand{\tightlist}{%
  \setlength{\itemsep}{0pt}\setlength{\parskip}{0pt}}
\setcounter{secnumdepth}{-\maxdimen} % remove section numbering
\usepackage{bookmark}
\IfFileExists{xurl.sty}{\usepackage{xurl}}{} % add URL line breaks if available
\urlstyle{same}
\hypersetup{
  pdftitle={Formative Assessment Report: Global Obesity and Diabetes Analysis},
  pdfauthor={Pranit Chatterjee},
  hidelinks,
  pdfcreator={LaTeX via pandoc}}

\title{Formative Assessment Report: Global Obesity and Diabetes
Analysis}
\author{Pranit Chatterjee}
\date{2025-11-06}

\begin{document}
\maketitle

\section{1. Business Understanding}\label{business-understanding}

\subsection{Introduction}\label{introduction}

This report investigates the relationship between obesity and diabetes
prevalence across countries and over time, with a particular focus on
understanding gender differences. Using data from the Non-Communicable
Disease Risk Factor Collaboration (NCD-RisC), we explore how these two
critical health risk factors have evolved from 1990 to 2022.

\subsection{Research Questions}\label{research-questions}

The analysis addresses the following key research questions:

\begin{enumerate}
\def\labelenumi{\arabic{enumi}.}
\tightlist
\item
  How does obesity prevalence compare between India and the United
  Kingdom over time?
\item
  Is there a relationship between obesity and diabetes prevalence
  globally?
\item
  Do gender differences exist in the obesity-diabetes relationship?
\item
  What trends can we observe in these health risk factors from 1990 to
  2022?
\end{enumerate}

\subsection{Key Performance Indicators
(KPIs)}\label{key-performance-indicators-kpis}

\begin{itemize}
\tightlist
\item
  \textbf{Prevalence of obesity (BMI ≥ 30 kg/m²):} Percentage of
  population classified as obese
\item
  \textbf{Prevalence of diabetes (age 18+ years):} Age-standardized
  percentage with diabetes diagnosis
\item
  \textbf{Temporal trends:} Changes in prevalence over 1990--2022
\item
  \textbf{Gender comparisons:} Sex-stratified analysis for Men and Women
\item
  \textbf{Country comparisons:} Focus on India, United Kingdom, and
  global aggregates
\end{itemize}

\subsection{Motivation}\label{motivation}

Understanding obesity and diabetes trends is critical for public health
planning, policy development, and resource allocation. These conditions
are major risk factors for cardiovascular disease, cancers, and other
non-communicable diseases (NCDs). By examining trends and relationships,
we can inform prevention strategies and identify high-risk populations.

\begin{center}\rule{0.5\linewidth}{0.5pt}\end{center}

\section{2. Data Understanding}\label{data-understanding}

\subsection{Data Sources}\label{data-sources}

This analysis uses two datasets from the NCD-RisC collaboration:

\begin{enumerate}
\def\labelenumi{\arabic{enumi}.}
\tightlist
\item
  \textbf{BMI Dataset}
  (\texttt{NCD\_RisC\_Lancet\_2024\_BMI\_age\_standardised\_country.csv})

  \begin{itemize}
  \tightlist
  \item
    Contains age-standardized prevalence of BMI categories
  \item
    Covers years 1990--2022
  \item
    Stratified by sex (Men, Women)
  \item
    Country-level and global aggregates
  \end{itemize}
\item
  \textbf{Diabetes Dataset}
  (\texttt{NCD\_RisC\_Lancet\_2024\_Diabetes\_age\_standardised\_countries.csv})

  \begin{itemize}
  \tightlist
  \item
    Contains age-standardized diabetes prevalence (18+ years)
  \item
    Covers years 1990--2022
  \item
    Stratified by sex (Men, Women)
  \item
    Country-level and global aggregates
  \end{itemize}
\end{enumerate}

\subsection{Data Loading and Initial
Exploration}\label{data-loading-and-initial-exploration}

\end{document}
